%============================================================ 2. 관련 연구 (압축판)
%  원판의 5개 소절을 3개로 합치고, 계열 열거를 대표 문헌 한 줄로 줄였다.
\section{관련 연구}
\label{sec:related}

\subsection{자산 배분과 무인 위협 대응}

방어 자산을 표적에 배분하는 문제는 무기-표적 할당(WTA)으로 정식화되어 정확해법과
휴리스틱이 모두 연구되었고~\citep{ahuja2007wta}, 1:1 배정으로 축소하면 헝가리안
알고리즘~\citep{kuhn1955hungarian}으로 최소비용 매칭을 얻는다. 다만 이전 결정이 이미 막은
회랑과 남은 자산에 따라 달라지는 판단은 고정된 거리·각도 비용으로 표현되지 않는다. 본
연구는 배정 계층에 조합 최적화 대신 언어모델의 판단을 둔다.

무인 위협의 무력화 수단은 연성(전자전)·경성(물리 타격)·포획으로
구분되는데~\citep{wang2021cuas}, 유선 조종 표적의 확산~\citep{sutea2026fiber}이 연성의
범위를, 요격 비용의 비대칭~\citep{rumbaugh2024cost}이 경성의 지속성을 제약하면서 포획이
남는다. 해상에서는 다중 USV 협력 기동~\citep{liu2016usv}과 강화학습 협조
포위~\citep{zhang2025encircle}가 보고되었으나 포위는 방어정이 표적보다 빠르거나 대등함을
전제하므로, 더 느린 조건에서는 추격 대신 접근 경로의 선제 차단이 남는다.

\subsection{다중 에이전트 강화학습과 크레딧 배분}

중앙 집중 학습·분산 실행(CTDE) 계열은 팀 보상에서 개체별 기여를 분리하려고 가치망을
학습하며, 중앙 비평가~\citep{lowe2017maddpg}, 단조 가치 분해~\citep{rashid2018qmix},
반사실적 기준선~\citep{foerster2018coma}, PPO~\citep{schulman2017ppo}의 다중 에이전트
확장~\citep{yu2022mappo}이 그 예다. 그러나 공유 가치망의 추정 편향과 학습
불안정~\citep{haarnoja2018sac}은 여러 에이전트가 함께 학습하는 비정상 환경에서 더 커진다.
그룹 상대 정책 최적화(GRPO)~\citep{shao2024deepseekmath}는 같은 상태에서 뽑은 여러 후보의
보상을 서로 비교해 이득을 산정함으로써 가치망 자체를 없앤다. 본 연구는 GRPO를 다중 에이전트
제어로 옮기되, 후보 하나를 기하 휴리스틱으로 고정해 기준선으로 삼고 후보의 팀 보상을
비평가가 아니라 시뮬레이터 재진행으로 얻는다. 다중 에이전트 근위 정책 최적화(MAPPO)는
같은 픽셀 행동공간 위의
기준선으로 두어 표현과 갱신 규칙의 기여를 분리한다(\ref{sec:res_algo}절).

\subsection{격자 표현·픽셀 지목과 언어모델 계획}
% Fig.2 는 §3.1 보다 먼저 보여야 한다. 2단 figure* 는 만난 쪽의 다음 쪽 상단에 놓이므로
% §3 이 시작하는 쪽 상단에 오도록 그 앞 쪽의 이 자리에서 선언한다.
\begin{figure*}[!t]
\centering
\includegraphics[width=\textwidth]{fig2.png}
\caption{세 가지 공격 포메이션: 집중·양동·파상(좌$\to$우)}
\label{fig:formations}
\end{figure*}

모든 픽셀에 값을 매기는 조밀 예측 계열~\citep{long2015fcn, ronneberger2015unet}과 그 확장인
좌표 채널~\citep{liu2018coordconv}·조건 변조~\citep{perez2018film}는 모두 상황 지각에
쓰였고, 후보마다 점수를 매겨 하나를 지목하는 출력~\citep{vinyals2015pointer,
vaswani2017attention}은 주로 조합 최적화에 머물렀다. 본 연구는 둘을 잇는다. U-Net 특징
지도의 각 픽셀에 방어정별 질의로 점수를 매기고(\ref{sec:policy}절), 그 점수맵을 정책의
행동 분포로 삼아 픽셀 하나를 지목하는 것이 곧 행동이므로, 안전 제약을 마스크로 행동 분포에
직접 부과할 수 있다.

언어모델을 계획 주체로 쓰는 연구는 고수준 지시를 실행 가능한 단계로
분해하고~\citep{huang2022zeroshot}, 실행 가능성 추정을 결합해 불가능한 계획을
걸러내며~\citep{ahn2022saycan}, 구조적 출력~\citep{openai2026function}으로 하위 시스템이 파싱
가능한 형태를 보장한다. 본 연구도 판단을 실행 가능성으로 거르되, 언어모델에는 배정만 맡기고
실행 가능성은 기동 계층의 유효 마스크가 픽셀 수준에서 강제한다(\ref{sec:validmask}절).
